%! Systems Involving Quadratics — Unit 2, Lesson 2.6 — Exit Ticket (Answer Key)
\documentclass[10pt]{article}
\usepackage{algebra2-article}
\usepackage{algebra2-key}   % pulls in -boxes; do NOT also load -boxes
\newcommand{\TallMath}[1]{$\displaystyle #1\rule[-1.4em]{0pt}{3.2em}$}

\begin{document}

\pageheader{Unit 2, Lesson 2.6}{Exit Ticket --- Answer Key}

\vspace{0.12in}

No notes. Show your steps. A solution is an ordered pair --- find \emph{both} coordinates and check them.

\begin{enumerate}[leftmargin=1.8em, itemsep=14pt]
  \item \textbf{Solve the system by substitution.} Write each solution as an ordered pair.
        $y=x^2-x-1, \qquad y=x+2$
        \begin{work}
          x^2-x-1 &= x+2 \\
          x^2-2x-3 &= 0 \\
          (x-3)(x+1) &= 0 \\
          x &= 3,\ -1 \\
          y &= 5,\ 1 \quad\text{so } (3,5) \text{ and } (-1,1)
        \end{work}

  \item \textbf{Number of solutions --- no solving.} How many solutions does the system
        $\ y=x^2+2,\ \ y=2x\ $ have? Substitute, then use the discriminant to justify.
        \[ b^2-4ac=\ans{$-4$} \qquad\Rightarrow\qquad \ans{$0$ solutions (the line misses the parabola)} \]

  \item \textbf{SOL-style multiple choice.} After substitution, a linear--quadratic system collapses to
        $x^2-6x+9=0$. How many solutions does the system have?
        \begin{enumerate}[label=\Alph*., leftmargin=1.9em, itemsep=1pt]
          \item $0$ solutions
          \item exactly $1$ solution \textcolor{keyred}{\textbf{$\leftarrow$ correct}}
          \item exactly $2$ solutions
          \item infinitely many solutions
        \end{enumerate}
        \begin{itemize}[leftmargin=1.4em, nosep]
          \item[] $x^2-6x+9=(x-3)^2=0$ has the double root $x=3$ ($b^2-4ac=0$), so the line is \emph{tangent} to the parabola --- exactly one shared point.
        \end{itemize}
\end{enumerate}

\end{document}
